\chapter{How Can Security Issue Be Addressed?}

\section{Local Network: Security Policy}

\subsection{Principles}
There are some common security attributes that should be present in any
system that processes valuable personal or sensitive information.
System designs should include mechanisms to enforce the following
security attributes.


\begin{htmllist}
  \htmlitemmark{RedBall}

  \latexhtml{\item {\bf Identification and authentication of
    users}}{\item[Identification and authentication of users]}

    Each user of a computer system should have a unique identification on
    the system, such as an account number or other user identification
    code. There must also be a means of verifying that the individual
    claiming that identity (e.g., by typing in that identifying code at a
    terminal) is really the authorized individual and not an impostor. The
    most common means of authentication is by a secret password, known only
    to the authorized user. But one of the most reliable means of
    authentication is by a smart card, owned only by the authorized user.


  \latexhtml{\item {\bf Authorization capability enforcing the principle of least
    possible privilege}}{\item[Authorization capability enforcing the
    principle of least possible privilege]}
    
    Beyond ensuring that only authorized individuals can access the system,
    it is also necessary to limit the users access to information and
    transaction capabilities. Each person should be limited to only the
    information and transaction authority that is required by their job
    responsibilities. This concept, known as the principle of least
    possible privilege, is a long-standing control practice. There should
    be a way to easily assign each user just the specific access
    authorities needed, and I believe that smart card technology could do
    this job well.


    \latexhtml{\item {\bf Individual accountability}}{\item[Individual
      accountability]}
    
    From both a control and legal point of view, it is necessary to
    maintain records of the activities performed by each computer user. The
    requirements for automated audit trails should be developed when a
    system is designed. The information to be recorded depends on what is
    significant about each particular system. To be able to hold
    individuals accountable for their actions, there must be a positive
    means of uniquely identifying each computer user and a routinely
    maintained record of each user's activities.


    \latexhtml{\item {\bf Audit mechanisms}}{\item[Audit mechanisms]}
    
    Audit mechanisms detect unusual events and bring them to the attention
    of management. This commonly occurs by violation reporting or by an
    immediate warning to the computer system operator. The type of alarm
    generated depends on the seriousness of the event.\\
    
    A common technique to detect access attempts by unauthorized
    individuals is to count attempts. The security monitoring functions of
    the system can automatically keep track of unsuccessful attempts to
    gain access and generate an alarm if the attempts reach an unacceptable
    number.


    \latexhtml{\item {\bf Performance assurance}}{\item[Performance assurance]}
    
    A basic design consideration for any information system should be the
    ability to verify that the system is functioning as intended. Systems
    that are developed without such design considerations are often very
    difficult to independently audit or review, leading to the possibility
    of unintended results or inaccurate processing.


    \latexhtml{\item {\bf Recoverability}}{\item[Recoverability]}
    
    Because companies can potentially be heavily dependent on a computer
    system, an important design consideration is the ability to easily
    recover from troublesome events, whether minor problems or major
    disruptions of the system. From a design point of view, systems should
    be designed to easily recover from minor problems, and to be either
    transportable to another backup computer system or replaced by manual
    processes in case of major disruption or loss of computer facility.
\end{htmllist}

\subsection{Access decisions}
Once the automated system is ready to use, decisions must be made
regarding access to the system and the information it contains. For
example, many individuals require the ability to access and view data,
but not the ability to change or delete data. Even when computer
systems have been designed to provide the ability to narrowly designate
access authorities, a knowledgeable and responsible official must
actually make those access decisions. The care that is taken in this
process is a major determining factor of the level of security and
control present in the system. If sensitive data is being transmitted
over unprotected lines, it can be intercepted or passive eavesdropping
can occur. Encrypting the files will make the data unintelligible and
port protection devices will protect the files from unauthorized
access, if warranted.


\subsection{Systems development process}
All information systems software should be developed in a controlled
and systematic manner according to company standards. The quality and
efficiency of the data processed, and the possible reconfiguration of
the system can all be affected by an inadequate development process.
The risk of security exposures and vulnerabilities is greatly reduced
when the systems development process is itself controlled.


\subsection{Physical security}
Information cannot be appropriately protected unless the facilities
that house the equipment are properly protected from physical threats
and hazards. The major areas of concern are described below.


\begin{htmllist}
  \htmlitemmark{RedBall}

    \latexhtml{\item {\bf Environmental conditions}}{\item[Environmental
      conditions]}
    
    For many types of computer equipment, strict environmental conditions
    must be maintained. Manufacturer's specifications should be observed
    for temperature, humidity, and electrical power requirements.


    \latexhtml{\item {\bf Control of media}}{\item[Control of media]}
    
    The media upon which information is stored should be carefully
    controlled. Transportable media such as tapes and cartridges should be
    kept in secure locations, and accurate records kept of the location and
    disposition of each. In addition, media from an external source should
    be subject to a check-in process to ensure it is from an authorized
    source.


    \latexhtml{\item {\bf Control of physical hazards}}{\item[Control of physical
      hazards]}
    
    Each area should be surveyed for potential physical hazards. Fire and
    water are two of the most damaging forces with regard to computer
    systems. Opportunities for loss should be minimized by an effective
    fire detection and suppression mechanism, and planning reduces the
    danger of leaks or flooding. Other physical controls include reducing
    the visibility of the equipment and strictly limiting access to the
    area or equipment.


    \latexhtml{\item {\bf Contingency planning}}{\item[Contingency planning]}
    
    Although risks can be minimized, they cannot be eliminated. When
    reliance upon a computer facility or application is substantial, some
    type of contingency plan should be devised to allow critical systems to
    be recovered following a major disaster, such as a fire. There are a
    number of alternative approaches that should be evaluated to most
    cost-effectively meet the company's need for continuity of service.


    \latexhtml{\item {\bf Configuration management}}{\item[Configuration
      management]}
    
    Risk can be introduced through unofficial and unauthorized hardware or
    software. Another key component of information resource management is
    ensuring only authorized hardware and software are being utilized.
    There are several control issues to be addressed.\\
    
    Records of hardware/software inventories, configurations, and locations
    should be maintained and kept up-to-date.


    \latexhtml{\item {\bf Complying with terms of software
      licenses}}{\item[Complying with terms of software licenses]}
    
    Especially with microcomputer software, illegal copying and other uses
    in conflict with licensing agreements are concerns. The use of software
    subject to licensing agreements must be monitored to ensure it is used
    according to the terms of the agreement.


    \latexhtml{\item {\bf Protecting against malicious software and
      hardware}}{\item[Protecting against malicious software and hardware]}
    
    The recent occurrences of destructive computer ``viruses'' point to the
    need to ensure that companies do not allow unauthorized software to be
    introduced to their computer environments. Unauthorized hardware can
    also contain hidden vulnerabilities. Management should adopt a strong
    policy against unauthorized hardware/software, inform personnel about
    the risks and consequences of unauthorized additions to computer
    systems, and develop a monitoring process to detect violations of the
    policy.
\end{htmllist}

\subsection{Data security}
Appropriate security mechanisms have to be in place in order to allow
responsible officials to designate access to data according to
individual computer users' specific needs. Security mechanisms should
be sufficient to implement individual authentication of system users,
allow authorization to specific information and transaction
authorities, maintain audit trails as specified by the responsible
official, and encrypt sensitive files if required by user management.


\subsection{Monitoring and review}
A final aspect of information resource protection to be considered is
the need for ongoing management monitoring and reviewing. To be
effective, a security program must be a continuous effort. Ideally,
ongoing processes should be adapted to include information protection
checkpoints and reviews. Information resource protection should be a
key consideration in all major computer system initiatives.\\

Earlier, the need for system audit trails was discussed. Those audit
trails are useful only if management regularly reviews exception items
or unusual activities. Irregularities should be researched and action
taken when merited. Similarly, all information-related losses and
incidents should be investigated.\\

A positive benefit of an effective monitoring process is an increased
understanding of the degree of information-related risk in agency
operations. Without an ongoing feedback process, management may
unknowingly accept too much risk. Prudent decisions about trade-off
between efficiency and control can only be made with a clear
understanding of the degree of inherent risk. Every manager should ask
questions and periodically review operations to judge whether changes
in the environment have introduced new risk, and to ensure that
controls are working effectively.


\subsection{Personnel management}
Information security is more a people issue than a technical issue.
Personnel are a vital link in the protection of information resources,
as information is gathered by people, entered into information resource
systems by people, and ultimately used by people. Security issues
should be addressed with regard to:
\begin{itemize}
  \item People who use computer systems and store information in the
        course of their normal job responsibilities.
  \item People who design, program, test, and implement critical or
        sensitive systems.
  \item People who operate computer facilities that process critical or
        sensitive data.
\end{itemize}

\subsection{Personnel security}
From the point of hire, individuals who will have routine access to
sensitive information resources should be subject to special security
procedures. More extensive background or reference checks may be
appropriate for such positions, and security responsibilities should be
explicitly covered in employee orientations. Position descriptions and
performance evaluations should also explicitly reference unusual
responsibilities affecting the security of information resources.\\

Individuals in sensitive positions should be subject to job rotation,
and work flow should be designed in such a way as to provide as much
separation of sensitive functions as possible. Upon decision to
terminate or notice of resignation, expedited termination or rotation
to less sensitive duties for the remainder of employment is a
reasonable precaution.


\subsection{Training}
Most information resource security problems involve people. Problems
can usually be identified in their earliest stages by people who are
attuned to the importance of information protection issues. A strong
training program will yield large benefits in prevention and early
detection of problems and losses. To be most effective, training should
be tailored to the particular audience being addressed.\\

Most employees want to do the right thing, if company expectations are
clearly communicated. Internal policies can be enforced only if staff
have been made aware of their individual responsibilities. All
personnel who access company computer systems should be aware of their
responsibilities under company policy, as well as obligations under the
law. Disciplinary actions and legal penalties should be communicated.


\subsection{User-level security policy}
Ultimately, computer security is the user's responsibility. Each user
must be alert to possible breaches in security and adhere to the
security regulations that have been established within their company.
The security practices listed in Appendix 1 are not inclusive, but
rather designed to remind them and raise their awareness towards
securing their information resources.


\section{Access Security: Firewalls}

\subsection{Introduction}
Many companies connect to the Internet, guarded by ``firewalls''
designed to prevent unauthorized access to their private networks.
Despite this general goal, firewalls span a continuum between ease of
use and security. This section describes some of the considerations and
tradeoffs in designing firewalls. A vocabulary for firewalls and their
components is given in Appendix to provide a common ground for
discussion.\\

The rationale for installing a firewall is almost always to protect a
private network against intrusion. In most cases, the purpose of the
firewall is to prevent unauthorized users from accessing computing
resources on a private network, and often to prevent unnoticed and
unauthorized export of proprietary information. In some cases export of
information is not considered important, but for many corporations that
are connecting this is a major though possibly unreasoning concern.
Many organizations will want simply to address the problem by not
connecting to the Internet at all. This solution can be difficult to
implement. If the private network is loosely administered or
decentralized, a single enterprising individual with a high-speed
dialup modem can quickly effect an Internet SLIP (Serial Line Internet
Protocol) connection which can compromise the security of an entire
network.\\

Often it is safe to say that a firewall needs to be put in place for
the ``CYA'' (Cover Your Assets, a family publication) factor. Even
though an employee could compromise proprietary information by carrying
it offsite on a DAT (Digital Audio Tape) or floppy disk, the Internet
represents a tangible threat, populated with dangerous ``vandals.''
(The Vandals were a collection of tribes of roughneck barbarians who
sacked Rome in 455 and looted it of all its portable wealth. Some use
the term ``hackers'' to describe Internet snoopers, but ``vandals,''
``crackers,'' or ``jerks'' is more appropriate) It could very easily
cost a network manager his job if a break-in occurs via this route,
even if the damage is no more extensive than could have been inflicted
over a dialup line or by a disgruntled employee. Generally, for a
would-be Internet site, the technical difficulties of implementing a
firewall are greatly outweighed by the public relations problems of
``selling'' upper management on the idea. In summary, because Internet
services are so highly visible, they are much more likely to require
official oversight and justification.


\subsection{Design decisions}
In configuring a firewall, the major design decisions with respect to
security are often already dictated by corporate or organizational
policy; specifically, a decision must be made as to whether security is
more important than ease-of-use, or vice versa. There are two basic
approaches that summarize the conflict:
\begin{itemize}
   \item That which is not expressly permitted is prohibited.
   \item That which is not expressly prohibited is permitted.
\end{itemize}

The importance of this distinction cannot be overemphasized. In the
former case, the firewall must be designed to block everything, and
services must be enabled on a case-by-case basis only after a careful
assessment of need and risk. This tends to impact users directly, and
they may see the firewall as a hindrance. In the second case, the
systems administrator is placed in a reactive mode, having to predict
what kinds of actions the user population might take that would weaken
the security of the firewall, and preparing defenses against them. This
essentially pits the firewall administrator against the users in an
endless arms race that can become quite fierce. Users can generally
compromise the security of their login if they try or are not aware of
reasonable security precautions. If the user has an open access login
on the firewall system itself, a serious security breach can result.
The presence of user logins on the firewall system tends to magnify the
problem of maintaining the system's integrity. A second important
statement of policy is implicit in the ``that which is not expressly
permitted is prohibited'' stance. This stance is more ``fail safe,''
since it accepts that the administrator is ignorant of what TCP ports
are safe, or what holes may exist in the manufacturer's kernel or
applications. Since many vendors are slow to publicise security holes,
this is clearly a more conservative approach. It is an admission of the
fact that what you do not know can hurt you.


\subsection{Levels of threat}
There are several ways in which a firewall can fail or be compromised.
While none of them are good, some are decidedly worse than others.
Since the purpose of many firewalls is to block access, it is a clear
failure if someone finds a loophole through it which permits them to
probe systems in the private network. An even more severe situation
would result if someone managed to break into the firewall and
reconfigure it such that the entire private network is reachable by
anyone. For the sake of terminology, this type of attack will be
referred to as ``destroying'' a firewall, as opposed to a mere
``break-in.'' It is extremely difficult to quantify the damage that
might result from a firewall's destruction. An important measure of how
well a firewall resists threat is the information it gathers to help
determine the course of an attack. The absolute worst thing that could
happen is for a firewall to be completely compromised without any trace
of how the attack took place. The best thing that can happen is for a
firewall to detect an attack, and inform the administrator politely
that it is undergoing attack, but that the attack is going to fail.\\

One way to view the result of a firewall being compromised is to look
at things in terms of what can be roughly termed as ``zones of risk.''
In the case of a network that is directly connected to the Internet
without any firewall, the entire network is subject to attack. This
does not imply that the network is vulnerable to attack, but in a
situation where an entire network is within reach of an untrusted
network, it is necessary to ensure the security of every single host on
that network. Practical experience shows that this is difficult, since
tools like rlogin that permit user-customizable access control are
often exploited by vandals to gain access to multiple hosts, in a form
of ``island hopping'' attack. In the case of any typical firewall, the
zone of risk is often reduced to the firewall itself, or a selected
subset of hosts on the network, significantly reducing the network
manager's concerns with respect to direct attack. If a firewall is
broken into, the zone of risk often expands again, to include the
entire protected network. A vandal gaining access to a login on the
firewall can begin an island hopping attack into the private network,
using it as a base. In this situation, there is still some hope, since
the vandal may leave traces on the firewall, and may be detected. If
the firewall is completely destroyed the private network can undergo
attack from any external system and reconstructing the course of an
attack becomes nearly impossible.\\

In general, firewalls can be viewed in terms of reducing the zone of
risk to a single point of failure. In a sense, this seems like a bad
idea, since it amounts to putting all of one's eggs in a single basket,
but practical experience implies that at any given time, for a network
of non-trivial size, there are at least a few hosts that are vulnerable
to break-in by even an unskilled attacker. Many corporations have
formal host security policies that are designed to address these
weaknesses, but it is sheer foolishness to assume that publishing
policies will suffice. A firewall enhances host security by funneling
attackers through a narrow gap where there is a chance of catching or
detecting them first. The well-constructed medieval castle had multiple
walls and interlocking defense points for exactly the same reason.


\subsection{Firewalls and their components}

\begin{htmllist}
    \htmlitemmark{RedBall}

    \latexhtml{\item {\bf Overview}}{\item[Overview]}
    
    In discussing firewalls there is often confusion of terminology since
    firewalls all differ slightly in implementation if not in purpose.
    Various discussions on Usenet indicate that the term ``firewall'' is
    used to describe just about any inter-network security scheme. For the
    sake of simplifying discussion, some terminology is proposed, to
    provide a common ground:\\
    
    \begin{htmllist}
      \htmlitemmark{WhiteBall}
    
    \latexhtml{\item {\bf Screening Router}}{\item[Screening Router (see
    \ref{scr_rout})]}
  
    A screening router is a basic component of most firewalls. A screening
    router can be a commercial router or a host-based router with some kind
    of packet filtering capability.  Typical screening routers have the
    ability to block traffic between networks or specific hosts, on an IP
    port level. Some firewalls consist of nothing more than a screening
    router between a private network and
    the Internet.\\
    
    \latexhtml{\item {\bf Bastion Host}}{\item[Bastion Host (see \ref{dual_hom})]}
    
    Bastions are the highly fortified parts of a medieval castle; points
    that overlook critical areas of defense, usually having stronger walls,
    room for extra troops, and the occasional useful tub of boiling hot oil
    for discouraging attackers. A bastion host is a system identified by
    the firewall administrator as a critical strong point in the network's
    security. Generally, bastion hosts will have some degree of extra
    attention paid to their security,
    may undergo regular audits, and may have modified software.\\
    
    \latexhtml{\item {\bf Dual Homed Gateway}}{\item[Dual Homed Gateway (see
    \ref{dual_hom})]}
    
    Some firewalls are implemented without a screening router, by placing a
    system on both the private network and the Internet, and disabling
    TCP/IP forwarding.  Hosts on the private network can communicate with
    the gateway, as can hosts on the Internet, but direct traffic between
    the networks is blocked. A dual homed gateway is, by definition, a
    bastion host.\\

\begin{figure}[htbp]
  \centering\includegraphics[width=4in]{dualhomed.ps}
  \caption{A typical Dual Homed Gateway}
  \label{fig:dual_hom}
\end{figure}

    \latexhtml{\item {\bf Screened Host Gateway}}{\item[Screened Host Gateway
    (see \ref{scr_hst})]}
    
    Possibly the most common firewall configuration is a screened host
    gateway. This is implemented using a screening router and a bastion
    host. Usually, the bastion host is on the private network, and the
    screening router is configured such that the bastion host is the only
    system on the private network that is reachable from the Internet.
    Often the screening router is configured to block traffic to the
    bastion host on specific ports, permitting only a small number of
    services to communicate with it.\\

\begin{figure}[htbp]
  \centering\includegraphics[width=4in]{screenhost.ps}
  \caption{A typical Screened Host Gateway}
  \label{fig:scr_hst}
\end{figure}

    \latexhtml{\item {\bf Screened Subnet}}{\item[Screened Subnet (see
    \ref{scr_sub})]}
    
    In some firewall configurations, an isolated subnet is created,
    situated between the Internet and the private network.  Typically, this
    network is isolated using screening routers, which may implement
    varying levels of filtering. Generally, a screened subnet is configured
    such that both the Internet and the private network have access to
    hosts on the screened subnet, but traffic across the screened subnet is
    blocked.  Some configurations of screened subnets will have a bastion
    host on the screened network, either to support interactive
    terminal sessions or application level gateways.\\

\begin{figure}[htbp]
  \centering\includegraphics[width=4in]{screensub.ps}
  \caption{A typical Screened Subnet}
  \label{fig:scr_sub}
\end{figure}
    
    \latexhtml{\item {\bf Application Level Gateway (or ``Proxy
    Gateway'')}}{\item[Application Level Gateway (or ``Proxy Gateway'')]}
    
    Much of the software on the Internet works in a store-and-forward mode;
    mailers and Usenet news collect input, examine it, and forward it.
    Application level gateways are service-specific forwarders or
    reflectors, which usually operate in user mode rather than at a
    protocol level.  Generally, these forwarding services, when running on
    a firewall, are important to the security of the whole. The famous
    sendmail hole that was exploited by the Morris Internet worm is one
    example of the kinds of security problems an application level gateway
    can present. Other application level gateways are interactive, such as
    the FTP and telnet gateways run on the Digital Equipment Corporation
    firewalls. In general, the term ``application level gateway'' will be
    used to describe some kind of forwarding service that runs across a
    firewall, and is a potential security concern. In general, crucial
    application
    level gateways are run on some kind of bastion host.\\
    
    \latexhtml{\item {\bf Hybrid Gateways}}{\item[Hybrid Gateways (see
    \ref{hyb_gat})]}
    
    Hybrid gateways are the ``something else'' category in this list.
    Examples of such systems might be hosts connected to the Internet, but
    accessible only through serial lines connected to an ethernet terminal
    server on the private network. Such gateways might take advantage of
    multiple protocols, or tunneling one protocol over another. Routers
    might maintain and monitor the complete state of all TCP/IP
    connections, or somehow examine traffic to try to detect and prevent an
    attack. The AT\&T corporate
    firewall is a hybrid gateway combined with a bastion host.\\
    
    Taking the components described above, we can accurately describe most
    of the forms that firewalls take, and can make some general statements
    about the kinds of security problems each approach presents. Assuming
    that a firewall fulfills its basic purpose of helping protect the
    network, it is still important to examine each type of firewall with
    respect to:\\
    
    Damage control --- If the firewall is compromised or destroyed, to what
    kinds of threats does it leave the private network open?\\
    
    Zones of risk --- How large is the zone of risk during normal
    operation?  A measure of this is the number of hosts or routers that
    can be probed
    from the outside network.\\

    Failure mode --- If the firewall is broken into or destroyed, how easy
    is this to detect? In a post mortem, how much information is retained
    that can be used to diagnose the attack?\\

    Ease of use --- How much of an inconvenience is the firewall?\\
    
    Stance --- Is the basic design philosophy of the firewall ``That which
    is not expressly permitted is prohibited'' or is it ``That which is not
    expressly prohibited is permitted''?
  \end{htmllist}

  \htmlitemmark{RedBall}
    \latexhtml{\item {\bf Firewalls using screening routers}}{\item[Firewalls
      using screening routers]}
    \label{scr_rout}
    
    Many networks are firewalled using only a screening router between the
    private network and the Internet. This type of firewall is different
    from a screened host gateway in that usually there is direct
    communication permitted between multiple hosts on the private network,
    and multiple hosts on the Internet. The zone of risk is equal to the
    number of hosts on the private networks, and the number and type of
    services to which the screening router permits traffic. For each
    service provided via peer-to-peer connection the size of the zone of
    risk increases sharply. Eventually it is impossible to quantify. Damage
    control is difficult as well since the network administrator would need
    to regularly examine every host for traces of a break-in. If there is
    no regular audit one must hope to stumble on a clue such as a
    mismatched system accounting record.\\
    
    In the case of total destruction of the firewall, it tends to be very
    hard to trace or even to discover. If a commercial router (which does
    not maintain logging records) is used, and the router's administrative
    password is compromised, the entire private network can be laid open to
    attack very easily. Cases are known where commercial routers have been
    configured with erroneous screening rules, or have come up in some
    pass-through mode because of hardware or operator error. Generally,
    this configuration is a case of ``That which is not expressly
    prohibited is permitted'' as the ingenious user can fairly easily
    piggyback protocols to achieve a higher level of access than the
    administrator expects or wants.\\
    
    Screening routers are not the most secure solution, but they are
    popular since they permit fairly free Internet access from any point
    within the private network. Many consultants and network service
    providers offer screening routers in a ``firewall'' configuration. It
    is uncertain if the various trade-offs involved are clear to the
    customer; thus the use of a screening router to protect sensitive
    information or trade secrets would not be recommended, since screening
    routers are very permeable from the inside.


    \latexhtml{\item {\bf Dual homed gateways}}{\item[Dual homed gateways]}
    \label{dual_hom}
    
    An often used and easy to implement firewall is the dual homed gateway.
    Since it doesn't forward TCP/IP traffic, it acts as a complete block
    between the Internet and the private network. Its ease of use is
    determined by how the systems manager chooses to set up access; either
    by providing application gateways such as telnet forwarders or by
    giving users logins on the gateway host. If the former approach is
    taken, the stance of the firewall is clearly ``That which is not
    expressly permitted is prohibited''; users can only access Internet
    services for which there is an application gateway. If users are
    permitted logins, the firewall's security is then seriously weakened.
    During normal operation, the only zone of risk is the gateway host
    itself, since it is the only host that is reachable from the Internet.
    If there are user logins on the gateway host, and one of the users
    chooses a weak password or has their account otherwise compromised, the
    zone of risk expands to encompass the entire private network. From a
    standpoint of damage control, the administrator may be able to track
    the progress of an intruder, based on the access patterns of the
    compromised login, but a skillful vandal can make this quite difficult.
    If a dual homed gateway is configured without direct user access,
    damage control can be somewhat easier, since the very fact that someone
    has logged in to the gateway host becomes a noteworthy security event.
    Dual homed gateways have an advantage over screening routers from the
    standpoint that their system software is often easier to adapt to
    maintain system logs, hard copy logs, or remote logs. This can make a
    post-mortem easier for the gateway host itself, but may or may not help
    the network administrator identify what other hosts on the private
    network may have been compromised in an island-hopping attack.\\
    
    Attacking a dual homed gateway leaves the attacker a fairly large array
    of options. Since the attacker has what amounts to local network access
    if a login can be obtained, all the usual attacks that can be made over
    a local network are available. NFS-mounted file systems, weaknesses in
    .rhosts files, automatic software distribution systems, network backup
    programs and administrative shell scripts -- all may provide a toehold
    on systems on the internal network. Once a toehold is secured, it then
    provides a base from which to launch attacks back at the gateway
    itself. The weakest aspect of the dual homed gateway is its failure
    mode. If the firewalll is destroyed it is possible that a skillful
    attacker might re-enable routing and throw the entire private network
    open to attack. In the usual Unix-based dual homed gateway, TCP/IP
    routing is often disabled by modifying a kernel variable named
    ipforwarding; if systems privileges can be obtained or stolen on the
    gateway, this variable can be changed. Perhaps this seems far-fetched,
    but unless great care is paid to monitoring the software revision
    levels and configuration on the gateway host, it is not improbable that
    a vandal with a copy of the release notes for the operating system
    version and a login can compromise the system.


    \latexhtml{\item {\bf Screened host gateways}}{\item[Screened host gateways]}
    \label{scr_hst}
    
    Several articles have described screened host gateways, and how to
    construct them. Generally, the screened host gateway is very secure,
    while remaining fairly easy to implement. Typically, a bastion host is
    configured on the private network, with a screening router between the
    Internet and the private network, which only permits Internet access to
    the bastion host. Since the bastion host is on the private network,
    connectivity for local users is very good, and problems presented by
    exotic routing configurations do not present themselves. If the private
    network is, as many are, a virtual extended local area network (e.g.:
    no subnets or routing) the screened host gateway will work without
    requiring any changes to the local network, as long as the local
    network is using a legitimately assigned set of network addresses. The
    zone of risk of a screened host gateway is restricted to the bastion
    host, and the screening router, and the security stance of the screened
    host gateway is determined by the software running on that system. If
    an attacker gains login access to the bastion host, there is a fairly
    wide range of options for attacking the rest of the private network. In
    many ways, this approach is similar to the dual homed gateway, sharing
    similar failure modes and design considerations with respect to the
    software running on the bastion host.


    \latexhtml{\item {\bf Screened subnets}}{\item[Screened subnets]}
    \label{scr_sub}
    
    A screened subnet is usually configured with a bastion host as the sole
    point of access on the subnet. The zone of risk is small, consisting of
    that bastion host or hosts, and any screening routers that make up the
    connections between the screened subnet, the Internet, and the private
    network. The ease of use and basic stance of the screened subnet will
    vary, but generally a screened subnet is appealing only for firewalls
    that are taking advantage of routing to reinforce the existing
    screening. This approach forces all services through the firewall to be
    provided by application gateways, and places the stance strongly in the
    ``That which is not expressly permitted is prohibited'' category.\\
    
    If a screened subnet based firewall with inter-network routing blocked
    is attacked with an intent to destroy it, the attacker must reconfigure
    the routing on three networks, without disconnecting or locking himself
    out, and without the routing changes being noticed. No doubt this is
    possible, but it can be made very difficult by disabling network access
    to the screening routers, or by configuring the screening routers to
    only permit access from specific hosts on the private network. In this
    case, an attacker would need to break into the bastion host, then into
    one of the hosts on the private network, and then back out to the
    screening router --- and would have to do it without setting off any
    alarms.\\
    
    Another advantage of screened subnets is that they can be put in place
    in such a way that they hide any accidents of history that may linger
    on the private network. Many sites that would like to connect to the
    Internet are daunted by the prospect of re-addressing and re-subnetting
    existing networks. With a screened subnet with blocked inter-network
    routing, a private network can be connected to the Internet and changed
    gradually to new subnet and network addresses. In fact, this approach
    has been observed to significantly accelerate the adoption of new
    network addresses on loosely controlled private networks. Users will be
    more receptive to changing their host addresses if they can realize the
    benefits of Internet connectivity thereby, since hosts that are not
    correctly addressed cannot use the firewall properly. In most other
    respects, the screened subnet is very much dependent on the suite of
    software running on the bastion host. Screening a whole subnet provides
    functionality similar to the dual homed gateway or screened host
    gateway; it differs primarily in the extra level of complexity in
    routing and configuration of the screening routers.


    \latexhtml{\item {\bf Hybrid gateways}}{\item[Hybrid gateways]}
    \label{hyb_gat}
    
    ``Security through obscurity'' is not sufficient in and of itself, but
    there is no question that an unusual configuration, or one that is hard
    to understand, is likely to give an attacker pause, or to make them
    more likely to reveal themselves in the process of trying to figure out
    what they are facing. On the other hand there is a real advantage to
    having a security configuration that is easy to understand, and
    therefore easier to evaluate and maintain. Since the hybrid gateway is
    mentioned here in the category of ``something else'' no attempt will be
    made to describe the indescribable. Some hypothetical hybrids may serve
    to show how hybrid gateways might differ from and be similar to the
    other types.\\
    
    Let us postulate a hybrid gateway that consists of a box sitting on the
    Internet, which is capable of routing traffic, but also maintains a
    complete notion of the state of every TCP connection, how much data has
    gone across it, where it originated, and its destination. Presumably,
    connections can be filtered based on arbitrarily precise rules, such
    as: ``permit traffic between host A on the private network and all
    hosts on network B on the Internet via the telnet service if and only
    if the connection originated from host A between the hours of 9:00 am
    and 5:00 pm and log the traffic.'' This sounds terrific, providing
    arbitrary control with great ease of use, but some problems simply
    refuse to go away. Consider that someone wishing to circumvent the
    firewall, who broke into the private network via an unguarded modem,
    might very easily set up a service engine that was piggybacked over the
    telnet port. This is actually a fairly easy firewall to destroy.\\
    
    Another hybrid gateway might take advantage of various forms of
    protocol tunneling. Suppose the requirement is to connect to the
    Internet with very tight restrictions, but that a high degree of
    connectivity is required between the private network and an external
    network that is somewhat trusted (for example, a corporate R\&D
    department needs to be able to run X-windows applications on a
    supercomputer at another facility). The usual archetypal gateways
    discussed here could provide general purpose e-mail connectivity, but
    for secure point-to-point communications, an encrypted point-to-point
    virtual TCP/IP connection might be set up with the remote system, after
    users had authenticated themselves with a cryptographic smart card.
    This would be extremely secure, and might be made fairly easy to use,
    but has the disadvantage that the protocol driver needs to be added to
    every system that wants to share communication. It is hard to make any
    guesses about the failure mode of such a system, but the zone of risk
    is neatly limited to all the hosts which are running the tunneling
    protocol driver, and to which the individual user has smart card
    access. Some of this might be implemented in hardware or in the routers
    themselves. In the future, it is likely that the rapid growth of the
    Internet will fuel more development in this area, and we will see
    various hybrid gateways arise. The basic issues surrounding configuring
    a firewall will probably remain the same as the ones discussed here.


    \latexhtml{\item {\bf Other firewall-related tools}}{\item[Other
      firewall-related tools]}
    
    There is active research and development on tools to aggressively seek
    out and identify weaknesses in an entire network, or to detect the
    patterns that might indicate when an attack is in progress. These tools
    range from the simple checklist to complex ``expert systems'' with
    inference engines and elaborate rule bases. Many firewalls today run
    software that is designed to go forth and gather information relating
    to possible attacks and their origins, often using and abusing tools
    like finger and SNMP. Unless true artificial intelligence is developed,
    however, these tools cannot guard against an unknown form of attack,
    since they cannot possibly match the creativity of a network vandal.
    While often billed as being ``proactive'' they are in fact reactive,
    and generally will serve only to catch systems crackers armed with last
    year's bag of tricks. Catching the small fry is still worth doing, but
    it is likely that they are less of a threat than the fellow who is so
    eager to break into your network that he is doing research and
    development in new system cracking techniques.
\end{htmllist}

\subsection{Comments}
In dealing with firewalls, it is simply not reasonable to say that any
particular approach is best, since there are so many factors that
determine what the best firewall for a given situation may be. Cost,
corporate policy, existing network technology, staffing, and
intra-organizational politics may all easily outweigh the technical
considerations presented here.\\

There are a few observations worth making about firewalls at a very
general level. Firstly, a firewall is a leverage-increasing device from
a network management point of view. Rather than looking at it as ``all
eggs in one basket,'' it can also be viewed as a trustworthy basket,
and a single point from which a very important security system can be
controlled. The size of the zone of risk is crucial to the design; if
it is small, security can be maintained and controlled easily but if
security is compromised, the damage can be more severe. The ideal would
be to have such strong host-based security that a firewall would be
redundant. Systems administration costs and a hard dose of reality
prevents this ideal from being obtainable.\\

A second important aspect of firewall building is that it is not
something to undertake in a vacuum. Many sites are connected with a
simple firewall consisting of a screening router and nothing more
because someone told them that it was ``secure enough.'' In setting up
a firewall one must trade off time and money, security, and risk.
``Secure enough'' depends largely on what you risk losing. Finally, it
is important when approaching implementing a firewall to avoid the urge
to start from scratch. System security is a lot like pregnancy; one is
seldom only broken into a little bit, and it only takes a little
mistake or a moment of inattention to find oneself in a delicate
position. Leaning on the experiences of others, and learning from their
mistakes and successes is very important. Setting up a firewall is
definitely an area where having a wide background in experience to draw
upon is important. The vandals on the network have a wide background in
experience to draw upon as well, and a firewall administrator must
communicate with others, and must keep up to date on other
firewall-related happenings on the network. Static defenses do not work
unless they keep up with emerging tricks of the trade, or one's
firewall may be the next Maginot Line.


\section{Data Security: Cryptography}
\label{data_crypto}

\subsection{What is cryptography?}
According to a dictionary it is ``The art or process of writing in or
deciphering secret code.'' In terms of the topic at hand, it is the
transforming of information into a form unreadable by anyone without a
secret decryption key.


\subsection{What is a key?}
A key is a piece of data that can be used to encrypt or decrypt a given
message.\\

In a common form of cryptography, secret key cryptography, the sender
uses a secret key to encrypt some information. The sender must then get
both the secret key and the encrypted information to the receiver
securely. The problem with secret key cryptography, from the standpoint
of transactions over the Internet, is that anyone who can get both the
key and the encrypted information can decrypt the information. As we
noted above, it is fairly easy for an unauthorized person to intercept
my information on the Internet. If I send both an encrypted message and
the secret key that decrypts the message over the Internet, then I risk
having both intercepted and having the message decrypted by the
interceptor.


\subsection{What is public key cryptography (PKC)?}
Public key cryptography (PKC) is a system based on pairs of keys, a
public key and a private key, rather than on a single secret key. The
keys are related in such a way that information that is encrypted with
the public key can only be decrypted with the private key. This means
that, if I have a pair of keys, I can make my public key freely
available to anyone who wants it without worrying about who else might
get it. If anyone wants to send me a secret message, they can use my
public key to encrypt the message. Only I can decrypt the message,
using my private key, which I keep safe.\\

The action of encrypting information with PKC keys is significantly
slower than encrypting with a secret key, such as a DES (Data
Encryption Standard) key. So, in practice, a system that uses PKC first
generates a secret key according to a method such as DES. It uses the
secret key to encrypt the message, and it uses the PKC key to encrypt
the secret key. Then, it attaches the PKC-encrypted secret key to the
secret key-encrypted message.


\subsection{How does someone obtain a pair of keys?}
A specialized piece of computer hardware and/or software generates key
pairs on demand according to the method defined by the PKC standards.


\subsection{How is information integrity assured?}
When you put together your message to me and encrypt it with my public
key, you also have your computer calculate something called a ``message
digest'' and attach it to the message. The message digest is a small
piece of data that results from performing a particular mathematical
calculation on your message. It acts as the digital fingerprint of your
message --- the message digest of every message is different.\\

When I receive your message and decrypt it with my private key, I also
calculate the message digest of your message. Then, I compare the
digest I generated to the one you sent with your message. If they
match, then I know that I received exactly what you sent.


\subsection{What about authenticity?}
When you obtain my public key, from me or my company or from some other
organization, it comes with a ``certificate''. The certificate is a
piece of cryptographically secure information attached to my public key
that guarantees that the key you obtained is mine.


\subsection{Is the certificate trustworthy?}
Because it comes from an organization that you trust. If I hand you my
driver's license to identify myself, you accept it as valid
identification because you trust the organization that issued me the
license, namely the government of my state.\\

If I am a catalog sales merchant and you obtain my public key from some
place on the Internet, it might have both a certificate from my company
and a certificate from my company's bank attached to it. This is a way
of saying, ``Look, it's not just Acme Sales Corp. that says this public
key is from our designated representative, but First National Bank says
so, too.'' Presumably, for First National Bank to issue such a
certificate, I (or someone from my company) had to prove myself to be a
representative of Acme to the appropriate person at First National.\\

As you can see, there is a notion of a ``hierarchy of trust''
associated with public keys and certificates. At the top of the
hierarchy is some organization that we trust implicitly. With checks
and credit cards, we have to trust the banks that issue them. With
currency we have to trust the government that issued it.\\

A trusted organization that issues certificates on behalf of others is
called a ``certifying authority'' (CA).


\subsection{What about non-repudiatability of origin?}
Using my private key and the message I am sending, I have my computer
make a particular calculation that results in a piece of information
called a ``digital signature''. Since it involves my private key, the
signature can only have come from me. And since the calculation uses
the message I'm sending, the signature is only for this message. I then
have my computer attach the digital signature to the message, and I
send to you.\\

When you receive the message, your computer does a calculation that
uses the message, the digital signature and my public key. Since my
public and private keys are related, the calculation can verify 1) that
my private key was used to create the signature and 2) that the message
you received was used to create the signature.


\subsection{What about non-repudiatability of receipt?}
If you and I are representatives of companies conducting business with
each other, then we would both have PKC pairs that identify us. You
would have my public key and I would have yours. We might agree that,
once you receive a digitally signed message from me, you send me a
receipt. The receipt would be my digital signature from the message I
sent to you, with your digital signature attached to it. Using your
public key, I would verify your digital signature. Using my private key
I would verify my digital signature that you had sent back to me.\\

How can I tell when a message was originated?\\

I can attach a ``digital time-stamp'' to my message. First, I would
have my computer calculate my digital signature and attach it to the
original message. Next, I would have it calculate a message digest of
the message plus the signature. Next, I would send the message digest
(presumably over the Internet) to an external, trusted digital
time-stamping service (DTS). (Bellcore has created a subsidiary to
provide such a service.) The DTS would return a digital time-stamp
consisting of:
\begin{itemize}
  \item the message digest that I sent
  \item the date and time it was received by the DTS, and
  \item the digital signature of the DTS encrypted with the DTS's private
    key.
\end{itemize}

The ``message'' that the DTS signs would be the combination of my
message digest (the digital fingerprint of my original message) and the
time-stamp. When my computer received the time-stamp from DTS it would
attach it to my message, add my digital signature and send you the
whole thing.\\

When you received my message, your computer would verify the time-stamp
using the public key of the DTS, and it would verify my signature using
my public key.


\section{Transaction Security}

\subsection{What are Kerberos, PGP, PEM, SSL, S-HTTP and IPSEC?}

\subsubsection{Kerberos}

\begin{htmllist}
  \htmlitemmark{RedBall}
  
    \latexhtml{\item {\bf What is Kerberos? What is it good for?}}
    {\item[What is Kerberos? What is it good for?]}

    Kerberos is a network authentication system for use on physically
    insecure networks, based on the key distribution model presented by
    Needham and Schroeder (in ``Using Encryption for Authentication in
    Large Networks of Computers'', Communications of the ACM, Vol. 21(12),
    pp.  993-999, December, 1978). It allows entities communicating over
    networks to prove their identity to each other while preventing
    eavesdropping or replay attacks. It also provides for data stream
    integrity (detection of modification) and secrecy (preventing
    unauthorized reading) using cryptography systems such as DES.\\
    
    Kerberos works by providing principals (users or services) with tickets
    that they can use to identify themselves to other principals and secret
    cryptographic keys for secure communication with other principals. A
    ticket is a sequence of a few hundred bytes. These tickets can then be
    embedded in virtually any other network protocol, thereby allowing the
    process implementing that protocol to be sure about the identity of the
    principals involved.\\
    
    Practically speaking, Kerberos is mostly used in application-level
    protocols (ISO model level 7), such as Telnet or FTP, to provide user
    to host security. It is also used, though less frequently, as the
    implicit authentication system of data stream (such as {\tt
      SOCK\_STREAM}) or RPC mechanisms (ISO model level 6). It could also
    be used at a lower level for host to host security, in protocols like
    IP, UDP, or TCP (ISO model levels 3 and 4), although such
    implementations are currently
    rare, if they exist at all.\\
  
    There are several different versions and distributions of Kerberos.
    Most of them are based on an MIT distributions in one form or another,
    but the lineage is not always simple. Some of the distributions are
    freely available, some are stand-alone commercial products, and others
    are part of a larger free or commercial systems.\\
  
    Several vendors currently sell relevant Kerberos material, and to name
    a few: CyberSAFE, Cygnus Support, Digital Equipment Corporation, Emulex
    Network Systems, OpenVision Technologies, TGV.\\

    \latexhtml{\item {\bf What is Bones? What is it for?}}
    {\item[What is Bones? What is it for?]}

    Bones is a system that provides the Kerberos API without using
    encryption and without providing any form of security whatsoever. It is
    a fake that allows the use of software that expects Kerberos to be
    present when it cannot be.\\
    
    Why does it exist? Kerberos is a network security system which relies
    on cryptographic methods for its security. Since Kerberos' encryption
    system, DES, is not exportable, Kerberos itself cannot be exported or
    used outside of the United States in its original form.\\
    
    As a partial solution to this problem, the Kerberos source code was
    modified by the addition of {\tt \#ifdef NOENCRYPTION} around all calls
    to DES functions. Compiling this version with the symbol NOENCRYPTION
    defined results in a system that looks like Kerberos from an
    application's point of view but that does not require DES libraries
    (and, as a result, does not speak the real Kerberos protocol and does
    not provide
    any security).\\
    
    The final piece in this puzzle is a program called ``piranha'' which
    takes the Kerberos sources and removes all of the calls to the
    encryption routines, replacing it with the code which was under {\tt
      \#ifdef NOENCRYPTION}, producing the system known as Bones. Bones has
    the property that there is absolutely no question about whether or not
    it is legal to transport its sources across national boundaries, since
    it neither has any encryption routines nor any calls to encryption
    routines.\\
    
    {\tt \#ifdef NOENCRYPTION} was not documented, and it was only intended
    to be used in the above manner. Someone who tries compiling Kerberos
    with that {\tt \#define} has in some sense ``voided the warranty'', and
    will get something which is both not secure and not Kerberos.
\end{htmllist}


\subsubsection{Pretty Good Privacy (PGP)}
\begin{htmllist}
  \htmlitemmark{RedBall}

    \latexhtml{\item {\bf A little history}}
    {\item[A little history]}
    
    PGP is a remarkable phenomenon. Largely the effort of a single person,
    Philip Zimmermann, it provides a confidentiality and authentication
    service that can be used for e-mail and file-storage applications. In
    essence, what Zimmermann has done is the following:

    \begin{itemize}
      \item Selected the best available cryptographic algorithms as building
        blocks.
      \item Integrated the algorithms into a general-purpose application
        that is independent of operating systems and processors and that is
        based on a small set of easy-to-use commands.
      \item Made the package and its documentation, including the
        sourcecode, freely available via the Internet, bulletin boards, and
        commercial networks such as CompuServe.
      \item Entered into an agreement with a company (ViaCrypt) to provide
        a fully compatible, low-cost commercial version of PGP.
    \end{itemize}

    From its beginnings just a few years ago, PGP has grown explosively and is
    now widely used. A number of reasons can be cited for this growth:
    
  \begin{itemize}
    \item It is available for free worldwide in versions that run on a
      variety of platforms, including DOS, Windows, Unix, and Macintosh. And
      the commercial version satisfies users who want a product that comes
      with vendor support.
    \item It is based on algorithms that have survived extensive public
      review and are considered extremely secure (RSA for public-key
      encryption, IDEA for conventional encryption, and MD5 for hash, or
      summary, coding).
    \item It has a wide range of applicability and can be used by
      corporations that want to enforce a standardized scheme for
      encrypting files and messages as well as by individuals who wish to
      communicate securely with people worldwide over the Internet and
      other networks.
    \item It was not developed by, nor is it controlled by, any
      governmental or standards organization. For people with an
      instinctive distrust of ``the establishment'', this makes PGP
      attractive.
  \end{itemize}

  Some people have hesitated to use PGP for several reasons. First, the
  developer lives and works in the United States, and because PGP is a
  cryptographic product, there is concern about export controls for
  organizations that wish to interoperate between U.S. and non-U.S.  sites.
  However, some recent versions of PGP have initially seen the light of day
  outside the United States, which would seem to make U.S.
  export controls irrelevant.\\
  
  Second, the freeware version of PGP was developed and disseminated
  without a license for the use of RSA, which is required within the United
  States but not outside. For people concerned with this issue, the
  commercial version of PGP has an RSA license. In addition, MIT has
  issued a new freeware version PGP with a licensed version of RSA.\\
  
  \latexhtml{\item {\bf What PGP does}}{\item[What PGP does]}
    
    PGP provides two services: encryption and digital signatures.
    Encryption allows a user to encode a file for storage locally or for
    transmission as an e-mail message. The local storage option is handy if
    you are worried about other people having access to files on your
    machine. The e-mail option enables PGP to be used for private exchanges
    over a network. PGP encrypts the entire contents of the message in such
    a way that only the intended recipient can decode and read the message.
    Anyone else who attempts to capture or copy the message en route will
    receive meaningless garble.\\
    
    The digital signature service allows a user to `sign' a document before
    transmission in such a way that anyone can verify that the signature is
    genuine and belongs with a particular document. If someone alters the
    message or substitutes a different message, the signature will no
    longer be valid. And any recipient can verify that the message has been
    signed by its true creator and not an imposter.\\
    
    PGP's confidentiality and encoding services use the most popular
    public-key encryption scheme, known as RSA. All public-key encryption
    systems make use of an encoding and decoding algorithm and a related
    pair of keys. The input to the encryption algorithm is the text to be
    encrypted --- known as plaintext --- and a key. The algorithm takes the
    input and produces scrambled output known as ciphertext. To use the
    decryption algorithm, you input the ciphertext plus the key that
    matches the one used for encryption,
    and the original plaintext is produced as output.\\
    
    The two keys used in any public-key encryption scheme, including RSA,
    are called the public key and private key. The public key, as the name
    suggests, is made public. The idea is to make your public key available
    to people with whom you correspond. You keep your private key secure,
    and it should be known only to you. These two keys can be used to
    provide confidentiality and encryption.\\
    
    Here's how it works. Suppose Alice wants to correspond with Bob. If
    Alice prepares a message and encrypts it with Bob's public key, only
    Bob can decrypt the message using his private key. If Alice prepares a
    message and encrypts it with her private key, then anyone, including
    Bob, can decrypt the message. But only Alice could have encrypted the
    message, therefore the encrypted message is, in effect, signed by Alice.\\
    
    It turns out that RSA, and all other known public-key algorithms, are
    time-consuming and inefficient. Therefore PGP, like most other
    encryption applications, does not use RSA directly to provide
    confidentiality and digital signatures.\\
    
    For confidentiality, PGP encrypts messages with an efficient single-key
    or conventional encryption algorithm known as IDEA. It then uses RSA to
    encrypt, with the receiver's public key, the IDEA key used to encrypt
    the message. The receiver can use RSA to recover the IDEA key and use
    that key to recover the message.\\
    
    For digital signatures, PGP uses an efficient algorithm known as MD5 to
    produce a summary code, or hash code, of the message that is, for all
    practical purposes, unique to that message. PGP then uses RSA to
    encrypt the hash code with the sender's private key. The receiver can
    use RSA to recover the hash code and verify that it is the correct hash
    code for the message. If it is correct, then only the alleged sender
    could have prepared the encrypted hash code.\\

  \latexhtml{\item {\bf Trusting public keys}}{\item[Trusting public keys]}
    
    The most difficult aspect of using PGP, or any public-key application,
    is getting your hands on the public keys of the people with whom you
    wish to correspond. You must make sure you have the true public key of
    each individual in your electronic Rolodex. Suppose I create a linked
    pair of public and private keys and send you the public key, declaring
    that I am Elvis. How do you know I am the real Elvis and not an
    impostor? If I am an impostor, I could send you signed messages and you
    would be sure they were from Elvis. If you send an encrypted message to
    Elvis, I can capture the message and recover the plain text.\\
    
    PGP provides a number of tools and recommended procedures for obtaining
    public keys in trusted ways. One handy tool is the public-key
    fingerprint, which is nothing more than a string of printable
    characters based on the MD5 hash code of the key. For all practical
    purposes, the fingerprint of a key is unique. So, if Alice knows Bob's
    voice, Bob could send his public key to Alice via e-mail. Alice then
    could generate the fingerprint of that key, call Bob, and have Bob read
    the fingerprint over the phone to make sure there is a match.\\
    
    Once you have a few trusted keys, you can make use of PGP's signature
    capability. If you have Bob's public key and you trust Bob to provide
    you with other public keys of other persons, Bob can send you John's
    key signed by Bob. That is, Bob takes John's public key and feeds it
    through the signature mechanism of PGP. Alice can use Bob's public key
    to ensure that John's key was provided by Bob and that the key has not
    been altered.\\
    
    There also are a number of servers on the Internet that are public-key
    repositories. Most of keys are signed by one or more people. You can
    obtain someone's public key from the server and if you trust the
    signatories to the key, you can have faith that it is genuine. These
    public-key servers do not authenticate the keys; they merely serve as
    repositories.\\
    
    One public-key directory that does attempt to provide authenticated PGP
    keys is SLED (Stable Large E-mail Database). The public keys in the
    directory are signed by SLED, indicating that the user's authenticity
    has been verified.
\end{htmllist}

\subsubsection{Privacy Enhanced Mail (PEM)}

\begin{htmllist}
  \htmlitemmark{RedBall}
  \latexhtml{\item {\bf What is PEM?}}{\item[What is PEM?]}
    
    PEM is an Internet standard for providing security services to
    electronic mail. It uses cryptographic techniques to provide message
    integrity checking, originator authentication, and confidentiality. It
    lets you know that a message hasn't been changed, who it's from, and,
    optionally, allows you to keep it secret from all but the intended
    recipients.\\
    
    There is a set of Proposed Standard RFCs (Internet standards documents)
    that specify PEM. The four documents are RFCs 1421 (obsoletes 1113),
    1422 (obsoletes 1114), 1423 (obsoletes 1115), and 1424 (new). The
    integration of MIME (Multipurpose Internet Mail Extensions) and PEM
    will either produce a new RFC for message format or cause RFC 1421 to
    be updated.\\
    
    An overview of PEM was presented in the August 1993 issue (Volume 38,
    Number 8) of ``Communications of the ACM'' in an article entitled
    ``Internet Privacy Enhanced Mail'' by Stephen T. Kent.\\

  \latexhtml{\item {\bf What is TIS/PEM?}}{\item[What is TIS/PEM?]}
    
    Trusted Information Systems (TIS), under ARPA sponsorship and in
    cooperation with RSA Data Security Incorporated (RSADSI), has released
    a reference implementation of Privacy Enhanced Mail (TIS/PEM) to the
    Internet community.\\
    
    TIS/PEM is a Unix-based implementation (The authors are currently
    pursuing porting it to other operating systems) that has been
    integrated with Rand MH 6.8.3 and is easily integrated into other mail
    user agents. TIS/PEM is distributed in source form. It is openly
    available within the United States and Canada for non-commercial use
    (not for resale), but the export from the United States of the
    cryptography used in TIS/PEM is controlled by the United States
    government.\\

  \latexhtml{\item {\bf What is RIPEM?}}{\item[What is RIPEM?]}
    
    Riordan's Internet Privacy Enhanced Mail (RIPEM) is an implementation
    of PEM. RIPEM allows electronic mail to have the four security
    facilities provided by PEM: disclosure protection (optional),
    originator authenticity, message integrity measures, and
    non-repudiation of origin (always). RIPEM is not really a complete
    implementation of PEM, because PEM specifies certificates for
    authenticating keys, which RIPEM does not handle at this time (however,
    their addition is planned). They have already been added to the
    Macintosh version.\\
    
    RIPEM was written primarily by Mark Riordan <mrr@scss3.cl.msu.edu>, but
    most of the code is now in the public domain, except for the RSA
    routines, which are a library called RSAREF licensed from RSA Data
    Security Inc. RIPEM has been ported to almost all platforms.\\

  \latexhtml{\item {\bf The certificates issue}}{\item[The certificates issue]}
    
    PEM users have the choice of using bare public keys or public keys
    embodied in X.509 certificates. If we use X.509 certificates to bind
    distinguished names to RSA public keys, it is not necessary to join the
    Internet certification hierarchy or otherwise pay to use PEM.
    
    TIS/PEM is capable of generating the bare keys and certificates that we
    may need. Joining the Internet certification hierarchy has the benefit
    of making it easier to verify others' mail and for them to verify ours.
    To join the Internet certification hierarchy, we must sign up our
    Certification Authority (CA) under a Policy-level Certification
    Authority (PCA).\\
    
    A distinguished name is a hierarchical, globally unique name used to
    identify something or someone.\\
    
    A Certification Authority (CA) vouches for the binding between users'
    distinguished names and RSA public keys within an organization or
    organizational unit. The CA's distinguished name is that of the
    organization or organizational unit and users' distinguished names are
    created by starting with the CA distinguished name and adding something
    to uniquely and unambiguously identify the user, like a common name.\\
    
    Policy-level Certification Authorities (PCAs) vouch for the binding
    between a CA's distinguished name and RSA public key. By joining a PCA,
    others can verify our PEM messages by following the certification path
    to the Internet PCA certificate without having to have retrieved our
    RSA public key using secure, out of band means. PCAs may also make CA
    Certificate Revocation Lists (CRLs) and certificates available and
    provide other services for its members.\\
    
    PCAs can be differentiated by the policy that they advertise. The
    policy includes the level of effort --- and associated assurance ---
    that a PCA uses to insure the correctness of the binding and the
    requirements they place on CAs which issue certificates under them.
    They can also be differentiated by the other services they offer and
    their price.\\
    
    Several PCAs exist as part of the Internet certification hierarchy,
    including PCAs at RSADSI and TIS, and more may come online in the near
    future. Individual PCAs will have their own price schedules.
\end{htmllist}

\subsubsection{Secure Sockets Layer (SSL)}

\begin{htmllist}
  \htmlitemmark{RedBall}
    
  \latexhtml{\item {\bf Introduction}}{\item[Introduction]}
    
    Netscape Communications has designed and specified a protocol for
    providing data security layered between application protocols (such as
    HTTP, Telnet, NNTP, or FTP) and TCP/IP. This security protocol, called
    Secure Sockets Layer (SSL), provides data encryption, server
    authentication, message integrity, and optional client authentication
    for a TCP/IP connection.\\
    
    SSL is an open, non-proprietary protocol, and has been submitted to the
    W3C working group on security for consideration as a standard security
    approach for World Wide Web browsers and servers on the Internet.
    Netscape Communications is working with the W3C and W3C members
    (including EIT) on developing and standardizing common, robust security
    mechanisms and protocols for the Internet. They will fully support such
    mechanisms and protocols as they become standardized.\\
    
    The SSL protocol specification is freely available and Netscape
    Communications encourage everyone to examine the protocol and/or
    implement it themselves.\\

  \latexhtml{\item {\bf How SSL works}}{\item[How SSL works]}
    
    SSL provides a security ``handshake'' that is used to initiate the
    TCP/IP connection. This handshake results in the client and server
    agreeing on the level of security they will use, and fulfills any
    authentication requirements for the connection. Thereafter, SSL's only
    role is to encrypt and decrypt the bytestream of the application
    protocol being used (for example, HTTP, NNTP, or Telnet). This means
    that all the information in both the HTTP request and the HTTP response
    are fully encrypted, including the URL the client is requesting, any
    submitted form contents (including things like credit card numbers),
    any HTTP access authorization information (usernames and passwords),
    and all the data returned from the server to the client.\\
    
    The Netscape Navigator includes embedded Certificate Authority (CA)
    keys for certain CAs, including Netscape Communications ` test CAs. As
    new CAs come online, they will embed their keys as well. These embedded
    keys allow the browser to verify the legitimacy of arbitrary servers.
    The Document Information dialog allow one to inspect both the identity
    of a given server as well as the identity of the CA that issued the
    server its certificate. SSL requires servers to have certificates
    issued by a Certificate Authority; and the Netscape Commerce Server
    includes a mechanism to easily acquire such a certificate.\\
    
    Because HTTP+SSL (or ``https'') and HTTP are different protocols and
    typically reside on different ports (443 and 80, respectively), the
    same server system can run both secure and insecure HTTP servers
    simultaneously. This means a server can provide some information to all
    users using no security, and other information only securely, if the
    webmaster so choose. For instance, ``storefront'' and merchandise
    catalog could be insecure, while ordering and payment documents and
    forms could be secure.
\end{htmllist}

\subsubsection{Secure HTTP (S-HTTP)}
Secure HTTP is an extension of HTTP that provides a variety of security
enhancements for the Web:

\begin{itemize}
  \item {\bf Authentication} lets clients make sure they are communicating
    with the correct server, and servers make sure they are communicating
    with the correct client. Authentication is performed using digital
    certificates issued by Certifying Authorities such as VeriSign.
  \item {\bf Encryption} makes data transferred over the network
    unintelligible to eavesdroppers.
  \item {\bf Digital Signatures} provide two benefits. First, they verify
    that data transferred over the network was not changed en route.
    Second, they provide non-repudiation: The receiver of data can prove to
    a third party that the sender really sent the data.
  \item {\bf Secure HTTP} is extremely flexible; it allows each application
    to configure the amount of security required.  A transmission from
    client to server or server to client can be signed, encrypted, both, or
    neither.
\end{itemize}

Secure NCSA Mosaic and Spry SafetyWEB are two browsers that support the
S-HTTP protocol.


\subsubsection{Comparison between SSL and S-HTTP}
Two comparisons of the S-HTTP and SSL protocols are possible: inherent
capabilities and current implementations. Keep in mind that network
security is a hot topic so developments continue at a rapid pace.\\

Because S-HTTP is an application-level protocol, it can provide
non-repudiation of individual requests or responses through digital
signatures. SSL is a lower-level protocol and does not have this
capability.\\

S-HTTP, being an application-level protocol, is able to work with
firewalls. SSL's transport-level encryption, in contrast, hides the
application-level protocol from firewalls. A firewall that relays an
SSL connection has no idea what data is being passed back and forth
over the connection.\\

S-HTTP is more flexible than SSL, in that an application can configure
the level of security it needs. Encryption and digital signatures can
be expensive to compute, so in principle this flexibility can allow a
server to handle more connections or respond more quickly. The other
side of this trade-off is that SSL, a lower-level protocol, may be
easier to optimize: It encrypts more bytes, but might balance this out
with a lower cost per byte. Since SSL has fewer options it should be
easier to set up and administer than S-HTTP.\\

Current SSL protocols and implementations are limited in a variety of
ways relating to the handling of digital certificates; S-HTTP
implementations are considerably more flexible. These SSL limitations
are not inherent in the nature of SSL and will likely be removed as
proves necessary.

\subsubsection{Internet Protocol Security (IPSEC)}

\begin{htmllist}
  \htmlitemmark{RedBall}

  \latexhtml{\item {\bf The IPSEC working group charter}}{\item[The IPSEC working
  group charter]}
    
    IPSEC is not a general security architecture for the Internet, but is
    instead focused on the IP layer. IPSEC working group will develop
    mechanisms to protect client protocols of IP. A security protocol in
    the network layer will be developed to provide cryptographic security
    services that will flexibly support combinations of authentication,
    integrity, access control, and confidentiality. The protocol formats
    for the IP Security Protocol (IPSP) will be independent of the
    cryptographic algorithm. The preliminary goals will specifically pursue
    host-to-host security followed by subnet-to-subnet and host-to-subnet
    topologies.\\
    
    Protocol and cryptographic techniques will also be developed to support
    the key management requirements of the network layer security. The key
    management will be specified as an application layer protocol that is
    independent of the lower layer security protocol. The protocol will
    initially support public key-based techniques. Flexibility in the
    protocol will allow eventual support of Key Distribution Center (KDC -
    such as Kerberos) and manual distribution approaches.\\

  \latexhtml{\item {\bf The IP authentication header (AH)}}{\item[The IP
  authentication header (AH)]}
    
    The IP Authentication Header is designed to provide integrity and
    authentication without confidentiality to IP datagrams. The lack of
    confidentiality ensures that implementations of the Authentication
    Header will be widely available on the Internet, even in locations
    where the export, import, or use of encryption to provide
    confidentiality is regulated. The Authentication Header supports
    security between two or more hosts implementing AH, between two or more
    gateways implementing AH, and between a host or gateway implementing AH
    and a set of hosts or gateways. A security gateway is a system which
    acts as the communications gateway between external untrusted systems
    and trusted hosts on their own subnetwork. It also provides security
    services for the trusted hosts when they communicate with the external
    untrusted systems. A trusted subnetwork contains hosts and routers that
    trust each other not to engage in active or passive attacks and trust
    that the underlying communications channel (e.g., an Ethernet) isn't
    being attacked.\\

  \latexhtml{\item {\bf The IP encapsulating security payload (ESP)}}{\item[The
  IP encapsulating security payload (ESP)]}
    
    The IP Encapsulating Security Payload (ESP) is designed to provide
    integrity, authentication, and confidentiality to IP datagrams. The ESP
    supports security between two or more hosts implementing ESP, between
    two or more gateways implementing ESP, and between a host or gateway
    implementing ESP and a set of hosts and/or gateways. A security gateway
    is a system which acts as the communications gateway between external
    untrusted systems and trusted hosts on their own subnetwork and
    provides security services for the trusted hosts when they communicate
    with external untrusted systems. A trusted subnetwork contains hosts
    and routers that trust each other not to engage in active or passive
    attacks and trust that the underlying communications channel (e.g., an
    Ethernet) isn't being attacked. Trusted systems always should be
    trustworthy, but in practice they often are not trustworthy.\\
    
    Gateway-to-gateway encryption is most valuable for building private
    virtual networks across an untrusted backbone such as the Internet. It
    does this by excluding outsiders. As such, it is often not a substitute
    for host-to-host encryption, and indeed the two can be and often should
    be used together.\\

  \latexhtml{\item {\bf For more information}}{\item[For more information]}
    
    You can see the IPSEC Working Group Charter at:
    \htmladdnormallink{ftp://ftp.ripe.net/ietf/ipsec/ipsec-charter.txt}
    {ftp://ftp.ripe.net/ietf/ipsec/ipsec-charter.txt}
    with all goals and milestones for this protocol, and/or see the Request
    For Comments (RFC) on the Security Architecture for the Internet
    Protocol at:
    \htmladdnormallink{http://www.cis.ohio-state.edu/htbin/rfc/rfc1825}
    {http://www.cis.ohio-state.edu/htbin/rfc/rfc1825}.
\end{htmllist}

\subsection{How can these protocols contribute to transaction security?}

\begin{htmllist}
  \htmlitemmark{RedBall}

    \latexhtml{\item {\bf Kerberos}}{\item[Kerberos]}
    
    It allows entities communicating over networks to prove their identity
    to each other while preventing eavesdropping or replay attacks. It also
    provides for data stream integrity (detection of modification) and
    secrecy (preventing unauthorized reading) using cryptography systems
    such as DES.\\
    
    It is important to realize that Kerberos is a one-trick pony. It
    provides for mutual authentication and secure communication between
    principals on an open network by manufacturing secret keys for any
    requestor and providing a mechanism for these secret keys to be safely
    propagated through the network. Kerberos does not, per se, provide for
    authorization or accounting, although applications that wish to can use
    their secret keys to perform those functions securely. Kerberos also
    does not provide password validation for individual workstations unless
    care is taken.\\
    
    It is also important to understand that using Kerberos on time-sharing
    machines greatly weakens its protections, since user's tickets are then
    only as secure as the `root' account (read: not very). Furthermore,
    dumb terminals and most X terminals do not understand the Kerberos
    protocol and as a result their cable connections remain insecure.


    \latexhtml{\item {\bf PEM}}{\item[PEM]}
    
    As far as I know, TIS/PEM and RIPEM are not used (or, at least, widely
    used) for securing electronic commerce transactions, but only for
    personal email messages. However, we may assume that some vendors will
    support PEM encrypted messages for their transaction processes if a
    significant number of customers desire so.


    \latexhtml{\item {\bf PGP}}{\item[PGP]}
    
    Some softwares developers have implemented a system by which NCSA
    Mosaic client and/or server call external PGP program which encrypt and
    decrypt their communications and thus provide secure communications
    between the server and the client, and ensure that a user is who he/she
    says he/she is.\\
    
    However, the most ``natural'' mean to securely communicate with PGP is
    by exchanging encrypted email messages, since this can be performed an
    all platforms.\\
    
    In both cases, the merchant can register his public key on a server
    such as SLED (http://www.four11.com) where the customer will be able to
    retrieve it. The merchant will have to provide the server with evidence
    (by another mean than the Internet) that he/she claims who he/she
    actually is, so that the customer will be able to trust his/her public
    key.\\
    
    Several electronic storefronts, such as NetMarket, support purchase
    orders or customer sensitive information encrypted with PGP.


    \latexhtml{\item {\bf SSL}}{\item[SSL]}
    
    We can enter our credit card number on a secure (https) browser (we may
    assume that other browsers than Netscape Navigator will support SSL in
    a near future) form and transmit the form over the Internet to a secure
    Commerce Server without risk of an intermediary obtaining our credit
    card information. The security features offered by SSL compliant
    technology protect commercial transactions, as well as all other
    communications, from misappropriation and fraud that could otherwise
    occur as information passes through Internet computers.\\
    
    Secure communications does not eliminate all of an Internet user's
    concerns. For example, the end-user must be willing to trust the server
    administrator with his or her credit card number before he or she enter
    into a commercial transaction. Security technology secures the routes
    of Internet communication; but does not protect from unreputable or
    careless people with whom we might choose to do business.\\
    
    The situation is analogous to telling someone our credit card number
    over the telephone. We may be secure in knowing that no one has
    overheard our conversation (privacy) and that the person on the line
    works for the company we wish to buy from (authentication), but we must
    also be willing to trust the person and the company.\\
    
    Server administrators must take additional precautions to prevent
    security breaches. To protect information, they must maintain physical
    security of their server computers and control access to software
    passwords and private keys.\\
    
    A lot of electronic storefronts, such as MarketNet, now support SSL
    mechanism.

    \latexhtml{\item {\bf S-HTTP}}{\item[S-HTTP]}
    
    S-HTTP specifies ways to add all the secure features needed in a
    transaction to the conversations that occur over the Internet between
    World Wide Web clients and servers. S-HTTP does not say how these
    elements will appear to users or what elements should be used with what
    transactions; that is up to the developers of the clients and servers
    and their customers.\\
    
    However, a few companies like OpenMarket and NCSA have already released
    Web secure servers and clients that support S-HTTP, Terisa proposes
    toolkits to develop such systems, and CyberCash has developed a payment
    system based on this protocol.


    \latexhtml{\item {\bf IPSEC}}{\item[IPSEC]}
    
    The primary objective of the IPSEC work is to ensure that IPv4 and IPv6
    will have solid cryptographic security mechanisms available to users
    who desire security. These mechanisms are designed to avoid adverse
    impacts on Internet users who do not employ these security mechanisms
    for their traffic. These mechanisms are intended to be
    algorithm-independent so that the cryptographic algorithms can be
    altered without affecting the other parts of the implementation. These
    security mechanisms should be useful in enforcing a variety of security
    policies.\\
    
    As far as I know, the IPSEC is not yet implemented: this protocol is in
    a stage of ``work in progress'', but it may be worth to keep looking at
    what it will become.
\end{htmllist}


\section{What a Smart Card is Intended for -- A Typical Example: the
  Authentication Issue}

\subsection{Smart card applications}

\subsubsection{Prepaid smart cards}

\begin{htmllist}
  \htmlitemmark{RedBall}

  \latexhtml{\item {\bf Why should we use prepaid smart cards?}}{\item[Why should
  we use prepaid smart cards?]}
    
  \begin{itemize}
    \item The telephone card is just one example of the smart card being
      used as a prepaid card.
    \item Prepaid cards generate a considerable advantage in extra cash flow.
    \item Prepaid cards are good advertising supports because users of
      prepaid cards for a particular service are a well defined population. A
      card user will usually carry the card on him for weeks or months and
      use it several times before the total value is used.
    \item Prepaid cards can be used in vending machines, car washes, gaming
      machines, automatic laundries, photocopiers and other areas where coins
      are used.
    \item Card operated payphones reduce theft, and increase convenience for
      users.
    \item Of the different card technologies such as smart cards, magnetic
      cards and optical cards, smart cards have been selected by the
      telecommunication industry because they give the best global solution
      in terms of equipment and card cost, security and reliability.
    \item The prepaid card is the simplest form of smart card.
    \item Telephones which are able to use smart cards open up a totally new
      range of possibilities, using more sophisticated cards for services
      like home billing from public payphones and memorized phone numbers.
  \end{itemize}

  \latexhtml{\item {\bf One example: making public telephones more
  convenient}}{\item[One example: making public telephones more convenient]}
    
    In 1984 all French payphones were coin operated. In 1985, after having
    rejected alternative technologies such as magnetic or optical cards on
    the grounds of security and global cost, France Telecom decided to
    start installing smart card payphones. In 1994 nearly all public
    payphones use prepaid smart cards. 100 million phone cards were sold in
    1993 and 115 million are expected to be sold in 1994. This is currently
    the largest smart card project in the world.\\
    
    Any person wishing to use a payphone can buy phone cards from tobacco
    stores, post offices and some newspaper kiosks. The cards hold either
    50 units (40 Francs = \$6.7) or 120 units (96 Francs = \$16). The
    duration of one unit depends on the distance and the time of day. A
    display on the telephone indicates how many units are left on the card.
    When no units are left a new card must be used.\\
    
    The use of smart cards for payphones has had dramatic effects on the
    image of the payphone. In 1985, 44,000 payphones were broken into. As
    these phones were systematically replaced by the new card payphones,
    the motive was removed and the number of phones broken into fell to a
    mere 900 in 1989 and about zero in 1994.\\
    
    Freed from the obligation of having the right small change at hand,
    people telephoned 50\% longer from card phones. As card operated public
    phones ceased to be a target for theft, they were less often out of
    order. A survey of payphone users showed that the satisfaction index
    for card phones is 20 points higher than for coin phones, despite
    having to pay in advance for calls from a card operated public phone.\\
    
    Now customers will also be given the possibility to store pre-recorded
    phone numbers or abbreviated dialing numbers and use their card to pay
    other services.\\
    
    Phonecards are used in two additional ways: as a novel advertising
    support and as a collectors item just like stamps.\\
    
    At the beginning of 1994, 66 operators in 60 countries are already
    using smart cards for payphones and, looking at the uptake trend, it is
    expected that over 100 countries will have adopted the smart card for
    payphones before the end of the century. The major smart card users
    apart from France are Germany and the Latin American countries. All
    systems are similar to the French system using a prepaid card. The
    world market for smart payphone cards is expected to leap from an
    estimated 250 million cards in 1993 to an estimated 1 billion cards in
    the year 2000.\\
    
    For payphones, smart cards have proved to be cost effective: they can
    bring new services and are readily accepted by users. The same benefits
    can be found in many other applications where coin operated machines
    have suffered from the cost and loss of service through vandalism, coin
    collection, and low reliability. Smart cards are already being used for
    gas metering, car parking, car washing, arcade games, and vending
    machines.\\
    
    If you put a coin into a machine today, you will most probably use a
    smart card for the same operation in five years time!\\
    
    Why are smart cards replacing coins in the world payphone market?\\
    
    There are several reasons:
    \begin{itemize}
      \item Reduced theft and maintenance cost.
      \item The electronics necessary to read a smart card are lower cost
        and more reliable than for any other card technology.
      \item Smart card payphones designed to accept simple prepaid smart
        cards allow the use of more sophisticated smart cards.
      \item Smart cards in payphones benefit users and operators.
    \end{itemize}
\end{htmllist}

\subsubsection{Electronic money}

\begin{htmllist}
  \htmlitemmark{RedBall}

  \latexhtml{\item {\bf Replacing coins and bank notes}}{\item[Replacing coins
  and bank notes]}
    
    In Singapore the CashCard, is starting to replace coins and bank notes
    for everyday purchases. The CashCard is a smart card which acts as an
    Electronic Purse, that is to say that it holds electronic money. As the
    national electronic purse program unfolds, it is predicted that most
    Singapore citizens will use the CashCard by 1997. They will then be
    able to pay with electronic cash for most everyday services, such as
    cinemas, car parks, tourist attractions, petrol, telephones, vending
    machines, sports, retail outlets, fast food, in fact virtually
    everywhere where cash is used today. The CashCard is available in
    several different initial values, S\$20, S\$50 and S\$100 (the
    equivalent of \$12, \$31 or \$62 in \$US). The card is not linked to a
    bank account, however the possibility of combining the electronic purse
    card with an ATM cash withdrawal card or credit card is quite possible
    and will most probably be used.\\

  \latexhtml{\item {\bf Why use electronic money?}}{\item[Why use electronic
  money?]}
    
    The main reasons for using electronic money are consumer convenience,
    reducing the cost and work associated with handling cash, and avoiding
    fraud. It becomes quite clear when examining these reasons that they
    are in fact quite closely linked.\\

  \latexhtml{\item {\bf Bringing more convenience for the
  consumer}}{\item[Bringing more convenience for the consumer]}
    
    Using a smart card to pay with electronic cash is faster and easier
    than using coins and bank notes; no more fumbling for that last little
    coin that you are certain that you saw in the bottom of your pocket
    (while 14 frustrated shoppers wait in a line behind you). Paying with
    electronic cash in most systems simply means checking that the amount
    on the till is the price you should be paying, and inserting the
    electronic purse in the terminal. Some systems may also have a button
    to push. Consumer research conducted in electronic purse pilot schemes
    show that consumer acceptance is very high.\\

  \latexhtml{\item {\bf Reducing the cost of handling cash}}{\item[Reducing the
  cost of handling cash]}
    
    The cash that has been counted and recounted every time that it enters
    or leaves the till for each sale during the day, is then counted and
    recounted and recorded to be taken to the bank, where it is counted
    again. With electronic money, there is no change to be given, the sales
    terminal automatically counts the taking and wires it directly to the
    bank.\\

  \latexhtml{\item {\bf Fighting fraud}}{\item[Fighting fraud]}
    
    The smart card, in association with a well designed security system,
    has been demonstrated to be an effective weapon against payment fraud.
    The electronic purse is also an effective protection against
    counterfeit coins and bank notes.\\

  \latexhtml{\item {\bf The electronic purse: a worldwide phenomenon}}{\item[The
  electronic purse: a worldwide phenomenon]}
    
    Singapore is not alone in moving towards an electronic purse system. At
    the start of 1994 over 20 projects had been identified spread over the
    5 continents, and to name a few, we can mention Mondex (which is still
    at an experimental stage in England and in the US, but which aims to
    become a worldwide electronic cash and payment systems standard), or
    Banksys --in Belgium-- and Interpay --in Holland-- Cash Cards (which
    claims to be the world's largest live electronic purse system). The
    technology is available, consumer acceptance has been demonstrated for
    card payment in the limited area of public payphones, and the benefits
    for consumers, merchants and banks are beginning to be well understood.
    By the year 2000, many working systems will be in place and perhaps in
    countries like Singapore, where an early adoption of the electronic
    purse has been made, electronic cash could well be more frequently used
    than today's coins and bank notes.
\end{htmllist}

\subsection{What is a smart card?}

\subsubsection{Presentation}
A smart card is a credit card sized plastic card with a special type of
integrated circuit embedded in it. The integrated circuit holds
information in electronic form and controls who uses this information
and how.\\

The type of information held, and the way this information is
controlled, gives rise to the infinite variety of smart card uses. The
plastic card has words and pictures printed on it which can be
understood by the card holder and people he meets. The integrated
circuit holds information in electronic form that can be easily,
securely and accurately accessed by all sorts of electronic data
processing equipment. This unique combination means that the smart card
is a bridge between man and machine.


\subsubsection{The plastic card}
The size of the card is determined by an international standard (ISO
7816). This standard also defines the physical characteristics of the
plastic, such as temperature tolerance and flexibility, the position of
the electrical contacts and their functions, and how the integrated
circuit talks to the outside world.\\

There are several types of plastic used for smart cards. The main types
are PVC and ABS. PVC can be embossed but is not recyclable. ABS cannot
be embossed but is recyclable. Within 5 years there are likely to be 1
billion smart cards used per year. So recycling cards will be important.


\subsubsection{Access control to the information}
Smart cards are smart because they control access to the information
they hold. Nobody can add units to a telephone card (not even the
telephone company). Only a doctor equipped with the correct card can
read a patient's medical record.\\

The basic mechanisms of access control are described below, and
illustrated using the examples mentioned in the first part of this
document.\\

There are two types of access control:
\begin{htmllist}
  \htmlitemmark{WhiteBall}
  \latexhtml{\item {\bf Who can access the information?}}{\item[Who can access
    the information?]}
  
    Everybody, the card holder or a specific third party?
    
  \latexhtml{\item {\bf How information can be accessed?}}{\item[How information
  can be accessed?]}
  
    Can it be read, added to or even modified or erased?
\end{htmllist}


\begin{htmllist}
  \htmlitemmark{RedBall}

  \latexhtml{\item {\bf Who can access the information?}}{\item[Who can access
  the information?]}
    
    Examples:
    \begin{itemize}
      \item I should be the only person able to use my credit card.
      \item Only retailers who are affiliated to the credit card scheme should
        be able to accept it.
      \item My regular doctor should be the only one who can add information to
        my medical record.
      \item Other doctors should be able to consult it.
      \item Emergency services should be able to know my blood group and
        address.
    \end{itemize}
    
    Some smart cards will only allow access to information if they are
    given the correct password. If you possess the password you can have
    access to the information. A password may be a number that the card
    holder has to remember, it may be hidden in another smart card that
    only authorized people possess, or it may just be inside the specific
    telephone or sales terminal used with the smart card. Most
    sophisticated smart cards have the ability to recognize several
    passwords.\\
    
    This means that in the case of a health card, the patient, the doctor,
    the hospital administration, and the health insurance organization
    could each have their own password giving them different access rights
    to sections of information held in a patient's health card.\\
    

    {\it Everybody}\\
    In some cases no password is used. So everybody can access the
    information. For example, in the health card, emergency medical
    information such as blood group and the patient's name are not
    protected by a password.\\
    
    {\it Card Holder Only}\\
    The most common form of password for card holders is a PIN (Personal
    Identification Number), a 4 or 5 digit number to be typed in at a key
    pad. The meaning of typing in a PIN is ``This card is mine, I give my
    permission to debit my account or access my medical file.'' In theory,
    if a smart card is lost or stolen and it is protected by a PIN, it
    cannot be used. If somebody tries to guess the PIN, most smart cards
    will just lock-up after 3 wrong attempts. PINs have two disadvantages,
    people can forget them, and it is possible to `steal ` a PIN by
    watching carefully when a card holder uses it. In the next few years
    better passwords for card holders which do not have these disadvantages
    are likely to appear.\\
    
    {\it Third Party Only}\\
    For example, when a doctor needs to consult a patient's health card, he
    needs a password, this password is hidden in his doctor's card. As only
    correctly trained doctors are issued with these cards, only they are
    able to consult patients' cards.\\

  \latexhtml{\item {\bf How can the information be accessed?}}{\item[How can the
  information be accessed?]}
    
    {\it Read only}\\
    Some information in the smart card can only be read. The information is
    fixed like a book, where information cannot be added, modified or
    erased. Some of the information in a prepaid telephone card is like
    this. This information contains a unique number for each card, the
    number of units in the new card, and an indication of who manufactured
    it. This information can be easily read, but cannot be modified in any
    way.\\
    
    {\it Add only}\\
    Some information in the smart card can only be added to but no erased,
    like engraving information in a stone. Information cannot be modified
    or erased, and information can only be added while there is room. If
    information can be added it can also be read. Units are used in a
    telephone card by adding information, like ticking off boxes. When all
    the boxes are ticked, no more units are left. Ticks can only be added
    and never erased, so nobody can refill the card.\\
    
    {\it Modify or erase}\\
    Information in the smart cards can be modified or erased like writing
    on a blackboard. This also implies that information can be added while
    there is room.\\
    
    {\it No access}\\
    Some information in a smart card can never be accessed. While this
    might appear rather useless, it is necessary for smart cards to be able
    to carry secrets the can never be revealed. The use of these secrets
    will be explained in ``When passwords are not enough.'' With a careful
    use of passwords, it is possible to avoid incorrect use of information,
    either by error or by intent. It is possible to summarize the
    information in a smart card, and the access control of this information
    (who and how) in a table. Take the example of a medical card in the
    table above. Information can be freely distributed in a smart card, and
    yet only accessed by the people who are authorized.\\

  \latexhtml{\item {\bf When passwords are not enough}}{\item[When passwords are
  not enough]}
    
    A smart card can restrict the use of information to an authorized
    person with a password. However if this information is then transmitted
    by radio or telephone additional protection is necessary.\\
    
    One form of protection is ciphering, which is like translating the
    information into some unknown foreign language. Imagine two Martians
    talking loudly in a London bar, everyone can hear, but only the
    Martians can understand and use the information.\\
    
    Some smart cards are capable of ciphering and deciphering (translating
    back to a form that can be easily understood), so that the information
    stored in them can be transmitted without compromising confidentiality.\\
    
    The ciphering scheme usually used is actually far more complex than the
    situation of our two Martian friends in London. Smart cards can cipher
    into billions and billions of ``foreign languages'', and choose a
    different language at random every time they start a conversation. This
    makes eaves-dropping virtually impossible.\\
    
    When a smart card and a computer system need to exchange confidential
    information over the telephone, they start by playing a complicated
    number game. The rules for this game are only known to genuine cards
    and genuine computers. It is impossible to play the game if the rules
    are not known, so any attempts to break into the system are easily
    detected and eliminated.\\
    
    If the game is correctly played, the genuine card knows that it is
    talking to a genuine computer and visa versa. This is known as
    authentication and some serious information exchange can now take place.\\
    
    During a successful authentication game one of the billions of possible
    ``foreign languages'' is selected, and only the players can know what
    it is. Information can then be transmitted in a language which is
    unknown to anyone else, even if they listened to the whole game.\\
    
    The ability of smart cards to enable totally secure communication is
    used in the GSM radio phone network. Even when the network covers the
    whole of Europe, only radio phones with a valid subscription card will
    be able to use it.\\

  \latexhtml{\item {\bf How smart is a smart card?}}{\item[How smart is a smart
  card?]}
    
    Some smart cards are smarter than others. The simplest cards, like the
    payphone cards, are smart enough to be virtually impossible to copy or
    falsify, but offer no protection in case of loss. If you lose your
    payphone card, anyone can use it. Other smart cards have one password
    to restrict its use to one person or machine, and the most
    sophisticated cards manage several passwords and can use authentication
    and ciphering techniques to combine total freedom with total security.\\
    
    The simplest smart cards can cost less than \$1 in large quantities,
    whereas the most sophisticated can cost in the range of \$6-10
    according to how much information they can hold.
\end{htmllist}

\subsection{How can your company use smart card systems?}

\subsubsection{Two-factor authentication process}
The major weakness in Kerberos based security is the reliance on
passwords for the authentication of the end user. Passwords need to be
reasonably easy to remember, or users will write them down, introducing
yet another security vulnerability. Yet easy-to-remember passwords use
only a small fraction of the key space, and are subject to dictionary
attacks.\\

Smart cards can dramatically reduce password exposure by supporting a
2-factor, 2-step authentication mechanism (See, for example, J. Kohl
and C. Neumann, ``The Kerberos Network Authentication Service'',
Internet Engineering Task Force, Networking Group, Internet Draft RFC
1510, Sep 1993 or Marjan Krajewski, Jr., ``Smart Card Augmentation of
Kerberos'', Proceedings of the Privacy and Security Research Group
Workshop on Network and Distributed System Security, February 11-12,
1993, San Diego, CA, pp.119-123.).\\

With passwords only, a user must know one piece of secret information
--- the valid password for a principal name --- in order to log in to a
company cell. There is one step involved in authentication --- supply a
login name and its password to the system.\\

With smart cards, a user must possess a unique physical token --- the
smart card --- as well as knowing one piece of secret information --- a
password associated with the card. Knowledge of the card password is of
no use without possession of the card itself. Possession of the card is
of no use without knowledge of the associated password.\\

With passwords only, there is one step to the authentication process ---
the user principal is authenticated to the computer. With smart cards,
there is an additional step involved --- before being authenticated to the
computer, the end user must be authenticated to the smart card. The first
step depends on knowing the correct password. The second step depends on
possessing the correct smart card --- the one containing the user's
long-term key.\\

Therefore it is hereby proposed using the smart card's secure file
system to store the user's long-term key. To access the key and log in
to the computer, a user must possess the card and must also know the
password required to access the key file on the card.\\

The details of this solution are contained in the following sections.


\subsubsection{Card requirements}
To achieve the additional security provided by this solution, cards
must meet certain minimum physical and functional requirements.

\begin{htmllist}
  \htmlitemmark{RedBall}

  \latexhtml{\item {\bf Physical access requirements}}{\item[Physical access
  requirements]}
    
    The card must conform to the accepted International Standards ISO
    7816-1, ISO 7816-2, and ISO 7816-3.\\
    
    The card must be tamper-resistant. Common tamper-resistance features
    include: chip enclosed in a hard, opaque, tamper-evident sealing
    coating over a passivation layer; circuitry to disable card output if
    the passivation layer is removed or if temperature, voltage, or
    frequency fluctuates outside the specified operating range.\\
    
    Cards with higher levels of tamper-resistance generally have a higher
    cost. The card issuer must evaluate with the card vendor the
    appropriate level of tamper-resistance for the issuer's application.\\
    
    Failure to provide adequate tamper-resistance makes the long-term key
    stored on the card vulnerable to physical attacks, reducing security
    when a card is lost or stolen.\\

  \latexhtml{\item {\bf Access control requirements}}{\item[Access control
  requirements]}
    
    The card file system must provide access controls that support the use
    of a password to protect access to the file containing the key.
    Passwords of up to 8 bytes in length must be supported. Without
    presentation to the card operating system of the password associated
    with the key file, there should be no access to the key stored in the
    card file system. Failure to conform to this specification negates the
    additional security that smart cards offer.\\
    
    The card must support detection of sequential invalid access attempts
    (a request for a protected card operation without supplying the correct
    password for the operation). Upon detecting N sequential invalid access
    attempts on any operation (where N may be configurable), the card must
    deny all further access to all operations except for possibly one
    ``unlock/erase'' operation. This one operation, if supported, must
    erase the key stored in the card file system before restoring normal
    access to the card. Failure to conform to this specification leaves a
    lost or stolen card vulnerable to off-line dictionary password-guessing
    attacks. Accounts associated with lost or stolen cards are subject to
    compromise until an administrator invalidates the account in the
    Security Registry.\\
    
    The card file system must support access controls that allow only the
    user to modify the password associated with the key file (with the
    possible exception described in the next paragraph).\\
    
    An interesting implementation would have a way of providing
    administrative access to unlock/erase a locked card (see above), set a
    new key on the card and in the Security Registry, set a new card access
    password, etc., but that is a vendor-issue not required for the basic
    functionality specified here.
\end{htmllist}

\subsection{The future}
Smart cards are just one particular way of using these special
integrated circuits. Historically it came about because credit cards
already existed, and the first smart card developments were for banking.\\

What is important about smart cards is that they are everyday objects
that people can carry in their wallets, purses or pockets, and that
they have the capacity to retain and protect information stored in
electronic form.\\

There are already other smart objects in use. Smart keys are being used
for pay-TV subscriptions. The smartness of smart cards comes from the
integrated circuit embedded in the plastic card. The same electronic
function could be performed by similar circuits embedded in other
everyday objects, key rings, watches, glasses, rings, or earrings.\\

The way that people give a password with a PIN is not yet ideal.
Biometric techniques will soon mean that a person can be reliably
identified by his hand, finger print, the color of his eyes or the
sound of his voice.\\

Soon it will be possible to hide electronic information in smart cards
and give authorization for this information to be used with a spoken
word or the touch of a hand.\\

The computer revolution has generated many advantages, but also
introduced much form filling, errors and impersonal service. Smart
cards along with other technologies such as biometrics, radio
telephones and home computing, will give us the advantages without the
disadvantages.\\

Smart cards are a relatively new technology that already effects the
everyday lives of millions of people. All the signs are that this is
just the beginning of something much larger that will influence the way
that we shop, see the doctor, use the telephone and enjoy leisure.\\

Modern society needs an enormous amount of information. Computers give
us the means to process this information. Smart cards give us a way of
humanizing the way we handle this information. Smart cards give every
individual the possibility to control how his personal information is
used.\\

Society will have to create laws which limit the uses of smart card
technology to those which benefit the individual, and increase the
quality of life.

